\documentclass{article}
\usepackage{../assignment_preamble}

\title{Homework 5}
\author{Ravi Kini}
\date{February 13, 2026}

\begin{document}

\maketitle

\problem{}
Let there be $n$ firms. Suppose the inverse demand function $P(Q)$ is:
\begin{align*}
    P(Q) = \begin{cases}
        a - Q & Q \leq a \\
        0 & Q > a \\
    \end{cases}
\end{align*}
and each firm has the cost function $c_i(q_i) = cq_i^2$, with $c < a$. Letting $q_{-i} = Q - q_i$, the profit function $\pi_i(q_i, q_{-i})$ is:
\begin{align*}
    \pi_i(q_i, q_{-i}) & = P(Q)q_i - cq_i^2 \\
    & = \max(a - q_i - q_{-i}, 0)q_i - cq_i^2
\end{align*}
Let us assume $a - q_i - q_{-i} \geq 0$. Taking the first derivative with respect to $q_i$:
\begin{align*}
    \frac{\partial \pi_i}{\partial q_i} = 0 & = (a - q_{-i}) - 2(1 + c)q_i \\
    q_i & = \frac{a - q_{-i}}{2(1 + c)}
\end{align*}
Now taking the second derivative:
\begin{align*}
    \frac{\partial^2 \pi_i}{\partial q_i^2} = - 2(1 + c)q_i < 0 \\
\end{align*}
The best response is then indeed $q_i^* = \frac{a - q_{-i}}{2(1 + c)}$, mking $(q_1^*, \ldots, q_n^*)$ a Nash equilibrium.

Suppose firms $i$ and $j$ have best responses $q_i^*$ and $q_j^*$, respectively, with $q_i^* \neq q_j^*$ and let ${Q'}^* = Q^* - q_i^* - q_j^*$. Then:
\begin{align*}
    2(1 + c)q_i^* & = a - {Q'}^* - q_j^*
    2(1 + c)q_j^* & = a - {Q'}^* - q_i^*
\end{align*}
It then follows that:
\begin{align*}
    2(1 + c)(q_i^* - q_j^*) & = q_i^* - q_j^* \\
    (1 + 2c)(q_i^* - q_j^*) & = 0
\end{align*}
For this to be true, $q_i^* = q_j^*$. This is a contradiction; therefore in a Nash equilibrium, all firms produce the same output $q^*$. That output is:
\begin{align*}
    q^* & = \frac{a - (n - 1)q^*}{2(1 + c)} \\
    q^* & = \frac{a}{1 + 2c + n}
\end{align*}
The equilibrium price is:
\begin{align*}
    P^* & = a - Q^* = a - n \cdot \frac{a}{1 + 2c + n} \\
    & = \frac{a(1 + 2c)}{1 + 2c + n}
\end{align*}
We see that $\lim_{n\to\infty} P^* = 0$, so the price decreases as the number of firms increases.
\end{document}